% !PRINTED 336
\weberpage{336}
Bezeichnen wir also mit $x_r$ Variable und setzen
\weq{7}{\overset{r}{\Sigma}\, x_r\, c^{(r)}_{h,k} = y_{h,k},}
so folgt:
\weq{8}{x_s = \overset{h,k}{\Sigma}\, a^{(s)}_{h,k}\, y_{h,k}.}
Jetzt bedeuten $u_r$, $v_r$, $t_r$ drei Systeme von Variablen und
\weq{9}{\begin{aligned}
\mu &= \Sigma\, \alpha_r\, u_r, \\
\nu &= \Sigma\, \beta_r\, v_r, \\
\lambda &= \Sigma\, \gamma_r\, t_r
\end{aligned}}
Basisformen von $\mathfrak{a}$, $\mathfrak{b}$, $\mathfrak{c}$, ferner $U$, $V$, $T$ die zu $\mathfrak{a}$, $\mathfrak{b}$, $\mathfrak{c}$ gehörigen Formen, in den Variablen $u,v,t$ geschrieben. Es ist dann
\weq{10}{\begin{aligned}
N(\mu) &= N(\mathfrak{a})\,U, \\
N(\nu) &= N(\mathfrak{b})\,V, \\
N(\lambda) &= N(\mathfrak{c})\,T.
\end{aligned}}
Es ergibt sich also aus (3) durch Multiplikation mit $u_h v_k$ und Summation nach (9):
\weq{11}{\mu\nu = \overset{r}{\Sigma}\, \gamma_r \overset{h,k}{\Sigma}\, c^{(r)}_{h,k}\, u_h v_k,}
also, wenn man
\weq{12}{t_r = \overset{h,k}{\Sigma}\, c^{(r)}_{h,k}\, u_h v_k}
setzt:
\weq{13}{\mu\nu = \lambda,}
und daraus, da $N(\mathfrak{a})N(\mathfrak{b}) = N(\mathfrak{c})$, ist nach (10)
\weq{14}{T = UV.}
Darin sind aber die $t_r$ nicht mehr unabhängige Variable, sondern sie gehen durch die \gesperrt{bilineare Substitution} (12) aus $u_h$, $v_k$ hervor.

5. \gesperrt{Die Form $T$ geht durch die bilineare Substitution (12) in das Produkt der beiden Formen $U, V$ über.}

Diese Substitution hat aber noch eine wesentliche Eigentümlichkeit. Wir nennen die durch (12) bestimmten Funktionen $t_r$ nach einem Primzahlmodul $p$ \gesperrt{linear unabhängig}, wenn aus der Kongruenz
\weq{15}{\Sigma\, x_r t_r \equiv 0 \pmod{p},}
in dem die $x_r$ ganze rationale Zahlen sind, folgt, daß diese Zahlen alle durch $p$ teilbar sind. Findet dies für jede beliebige Primzahl $p$ statt, so heißen die $t_r$ \gesperrt{schlechthin linear unabhängig}.
